\paragraph{4. Primitive twisting is an actual Euler system.}
Fix a constituent $\chi=\chi_h$ of primitive conductor
$D=|\Delta_h|=2^a d$, $d$ odd, and a basis $e_\chi$ of
$\mathbf Z_2(\chi)$. Put
$\Sigma=\operatorname{prime}(2Ndcd_s)$.
For every positive integer $m$ with $(m,cd_s)=1$, set
\[
 K_m=\mathbf Q(\zeta_m),\quad
 L_m=\mathbf Q(\zeta_{Dm}),\quad
 y_m=\operatorname{cor}_{L_m/K_m}
             (z_{Dm}\otimes e_\chi).                      \tag{K8}
\]
The character is trivial over $L_m$, so tensoring is an integral
map there. Formula (K7) places (K8) in the precise $j_*$ group
required by Kato~13.1 on its index set
$\operatorname{prime}(m)\cap\Sigma=\{2\}$. The larger definition is
permitted by (8.1.3), Sections~8.9--8.11 and Proposition~8.12:
$S_{Dm}=\operatorname{prime}(2NDm)$ is disjoint from
$\operatorname{prime}(cd_s)$, since $d\mid F$ and (K5) holds.
Each class uses the original $f,\xi,c,d_s$. We will also use this
same construction on the different index set required by
\cite[Section 0.1]{kato1999}, as checked below.

For $m\mid m'$, transitivity of corestriction reduces the norm
of (K8) to Proposition~8.12, p.~186. For a deck automorphism $g$,
equivariance gives
\[
 \operatorname{cor}((gz)\otimes e_\chi)
       =\chi(g)^{-1}g\operatorname{cor}(z\otimes e_\chi).
\]
Consequently the factor at a newly introduced good prime $q$ is
\[
 E_q(\chi)=1-\frac{a_q\chi(q)}q\sigma_q^{-1}
                         +\frac1q\sigma_q^{-2}.           \tag{K9}
\]
In Kato~13.1, pp.~224--225, the polynomial is
$P_q(X)=\det(1-\operatorname{Fr}_qX\mid T\chi)$ and the norm
factor is $P_q(q^{-1}\sigma_q^{-1})$, exactly (K9).
Repeated prime powers introduce no factor because Proposition~8.12
has a product over newly added support primes. Thus the subfamily
on Kato~13.1's index set is an Euler system of the stated lattice,
not a class inferred from proportionality of rank-one Iwasawa
modules. The required norm relations on the 1999 index set are
verified separately below.

